% Phantom Codes — supplementary materials master file (JAMIA supplement).
%
% Mirrors phantom_codes/main.tex's preamble (same biblatex setup, same
% packages) so prose conventions and citation rendering are identical
% between the main manuscript and the supplement. The main difference:
% this file \input{}'s build/supp_sections/*.tex (one per S1-Sn section)
% instead of build/sections/*.tex.
%
% Build:
%     cd paper && make phantom_codes_supp
%
% Output: paper/phantom_codes/build/supplementary.pdf (snapshotted to
%         paper/phantom_codes_supplementary.pdf via
%         `make snapshot-phantom_codes-supp`).
%
% Editing model:
%   - Prose edits happen in supp_sections/*.md
%   - This file only changes when you add/remove/reorder sections,
%     change document class, or restyle the title.

\documentclass[11pt]{article}

% Fonts & Unicode (XeLaTeX) — see main.tex for rationale.
\usepackage{fontspec}
\setmainfont{STIX Two Text}
\catcode`≥=\active \protected\def≥{\ensuremath{\geq}}
\catcode`≤=\active \protected\def≤{\ensuremath{\leq}}
\catcode`→=\active \protected\def→{\ensuremath{\rightarrow}}
\catcode`≈=\active \protected\def≈{\ensuremath{\approx}}

\usepackage[margin=1in]{geometry}
\usepackage{microtype}
\usepackage{parskip}

\usepackage{graphicx}
\usepackage{booktabs}
\usepackage{longtable}
\usepackage{array}
\usepackage{caption}
\captionsetup{font={small,it}, labelfont={small,bf,up}, labelsep=colon}
\usepackage{calc}

\usepackage{amsmath}
\usepackage{amssymb}

% biblatex with biber backend — matches main.tex; same Vancouver style
% so cross-references between main + supplement render consistently.
\usepackage[backend=biber, style=vancouver, sortlocale=en_US]{biblatex}
\addbibresource{../references.bib}

\usepackage[colorlinks=true,
            linkcolor=blue!60!black,
            citecolor=blue!60!black,
            urlcolor=blue!60!black]{hyperref}
\usepackage{xurl}

\usepackage{xcolor}
\usepackage{fancyvrb}
\usepackage{framed}
\usepackage{csquotes}

% Pandoc compatibility shims (same as main.tex)
\providecommand{\tightlist}{%
  \setlength{\itemsep}{0pt}\setlength{\parskip}{0pt}}
\providecommand{\pandocbounded}[1]{#1}

\title{%
  Supplementary Materials \\
  \large Phantom Codes: Hallucination in LLM-Based Clinical Concept Normalization}
\author{%
  Brian K. Fung \\[2pt]
  \small Independent Researcher \\[2pt]
  \small\href{mailto:brian@briankfung.com}{\texttt{brian@briankfung.com}}}
\date{\today}

\begin{document}

\maketitle

\noindent
This supplement accompanies the main manuscript. References cite the
shared bibliography (\texttt{paper/references.bib}); section labels use
the prefix S1--Sn for cross-reference from the main text.

\tableofcontents
\bigskip

% ─── Body ──────────────────────────────────────────────────────────────
% Each S* section is generated from supp_sections/*.md via pandoc; see
% Makefile. Add a new section by: (1) creating supp_sections/SN_name.md
% and (2) appending an \input{build/supp_sections/SN_name} line below.

\input{build/supp_sections/S1_prompt_templates}
\section{S2: Extended results}\label{s2-extended-results}

This supplement extends §3 of the main manuscript with the full
per-(model, mode) outcome distribution across all 28 configurations,
the complete Wilson 95\% CI tables for hallucination and
no\_prediction, the full top-k lift table, the per-bucket cost
decomposition, and an enumeration of every fabricated ICD-10-CM code
emitted in the headline run. All tables are auto-generated by
\texttt{phantom-codes\ report} from
\href{https://github.com/ufbfung/phantom-codes/blob/main/BENCHMARK.md}{\texttt{results/raw/headline\_n125\_20260504T040931Z.csv}}
into \href{https://github.com/ufbfung/phantom-codes/tree/main/results/summary/n125_run_v2/}{\texttt{results/summary/n125\_run\_v2/}};
this supplement reproduces them verbatim for reviewer auditability.

\subsection{S2.1: Full headline outcome distribution (28 configs × 4 modes)}\label{s2.1-full-headline-outcome-distribution-28-configs-4-modes}

Each row sums to n=125 across the six outcome buckets (counts and
percentages). The headline-relevant subset of these rows is in main
§3 Table 1; the full set is reproduced below for auditability.

\begin{longtable}[]{@{}
  >{\raggedright\arraybackslash}p{(\columnwidth - 16\tabcolsep) * \real{0.22}}
  >{\raggedright\arraybackslash}p{(\columnwidth - 16\tabcolsep) * \real{0.13}}
  >{\raggedleft\arraybackslash}p{(\columnwidth - 16\tabcolsep) * \real{0.05}}
  >{\raggedleft\arraybackslash}p{(\columnwidth - 16\tabcolsep) * \real{0.10}}
  >{\raggedleft\arraybackslash}p{(\columnwidth - 16\tabcolsep) * \real{0.10}}
  >{\raggedleft\arraybackslash}p{(\columnwidth - 16\tabcolsep) * \real{0.10}}
  >{\raggedleft\arraybackslash}p{(\columnwidth - 16\tabcolsep) * \real{0.10}}
  >{\raggedleft\arraybackslash}p{(\columnwidth - 16\tabcolsep) * \real{0.10}}
  >{\raggedleft\arraybackslash}p{(\columnwidth - 16\tabcolsep) * \real{0.10}}@{}}
\toprule()
\begin{minipage}[b]{\linewidth}\raggedright
Model
\end{minipage} & \begin{minipage}[b]{\linewidth}\raggedright
Mode
\end{minipage} & \begin{minipage}[b]{\linewidth}\raggedleft
n
\end{minipage} & \begin{minipage}[b]{\linewidth}\raggedleft
Exact
\end{minipage} & \begin{minipage}[b]{\linewidth}\raggedleft
Category
\end{minipage} & \begin{minipage}[b]{\linewidth}\raggedleft
Chapter
\end{minipage} & \begin{minipage}[b]{\linewidth}\raggedleft
OOD
\end{minipage} & \begin{minipage}[b]{\linewidth}\raggedleft
NoPred
\end{minipage} & \begin{minipage}[b]{\linewidth}\raggedleft
Halluc
\end{minipage} \\
\midrule()
\endhead
baseline:exact & D1\_full & 125 & 0 (0.0\%) & 0 (0.0\%) & 0 (0.0\%) & 0 (0.0\%) & 125 (100.0\%) & 0 (0.0\%) \\
baseline:exact & D2\_no\_code & 125 & 0 (0.0\%) & 0 (0.0\%) & 0 (0.0\%) & 0 (0.0\%) & 125 (100.0\%) & 0 (0.0\%) \\
baseline:exact & D3\_text\_only & 125 & 0 (0.0\%) & 0 (0.0\%) & 0 (0.0\%) & 0 (0.0\%) & 125 (100.0\%) & 0 (0.0\%) \\
baseline:exact & D4\_abbreviated & 125 & 0 (0.0\%) & 0 (0.0\%) & 0 (0.0\%) & 0 (0.0\%) & 125 (100.0\%) & 0 (0.0\%) \\
baseline:fuzzy & D1\_full & 125 & 67 (53.6\%) & 28 (22.4\%) & 0 (0.0\%) & 0 (0.0\%) & 0 (0.0\%) & 30 (24.0\%) \\
baseline:fuzzy & D2\_no\_code & 125 & 67 (53.6\%) & 28 (22.4\%) & 0 (0.0\%) & 0 (0.0\%) & 0 (0.0\%) & 30 (24.0\%) \\
baseline:fuzzy & D3\_text\_only & 125 & 67 (53.6\%) & 28 (22.4\%) & 0 (0.0\%) & 0 (0.0\%) & 0 (0.0\%) & 30 (24.0\%) \\
baseline:fuzzy & D4\_abbreviated & 125 & 25 (20.0\%) & 8 (6.4\%) & 30 (24.0\%) & 28 (22.4\%) & 0 (0.0\%) & 34 (27.2\%) \\
baseline:tfidf & D1\_full & 125 & 67 (53.6\%) & 15 (12.0\%) & 0 (0.0\%) & 0 (0.0\%) & 7 (5.6\%) & 36 (28.8\%) \\
baseline:tfidf & D2\_no\_code & 125 & 67 (53.6\%) & 15 (12.0\%) & 0 (0.0\%) & 0 (0.0\%) & 7 (5.6\%) & 36 (28.8\%) \\
baseline:tfidf & D3\_text\_only & 125 & 67 (53.6\%) & 15 (12.0\%) & 0 (0.0\%) & 0 (0.0\%) & 7 (5.6\%) & 36 (28.8\%) \\
baseline:tfidf & D4\_abbreviated & 125 & 25 (20.0\%) & 7 (5.6\%) & 16 (12.8\%) & 38 (30.4\%) & 0 (0.0\%) & 39 (31.2\%) \\
claude-haiku-4-5:constrained & D1\_full & 125 & 125 (100.0\%) & 0 (0.0\%) & 0 (0.0\%) & 0 (0.0\%) & 0 (0.0\%) & 0 (0.0\%) \\
claude-haiku-4-5:constrained & D2\_no\_code & 125 & 125 (100.0\%) & 0 (0.0\%) & 0 (0.0\%) & 0 (0.0\%) & 0 (0.0\%) & 0 (0.0\%) \\
claude-haiku-4-5:constrained & D3\_text\_only & 125 & 117 (93.6\%) & 8 (6.4\%) & 0 (0.0\%) & 0 (0.0\%) & 0 (0.0\%) & 0 (0.0\%) \\
claude-haiku-4-5:constrained & D4\_abbreviated & 125 & 117 (93.6\%) & 8 (6.4\%) & 0 (0.0\%) & 0 (0.0\%) & 0 (0.0\%) & 0 (0.0\%) \\
claude-haiku-4-5:rag & D1\_full & 125 & 125 (100.0\%) & 0 (0.0\%) & 0 (0.0\%) & 0 (0.0\%) & 0 (0.0\%) & 0 (0.0\%) \\
claude-haiku-4-5:rag & D2\_no\_code & 125 & 123 (98.4\%) & 2 (1.6\%) & 0 (0.0\%) & 0 (0.0\%) & 0 (0.0\%) & 0 (0.0\%) \\
claude-haiku-4-5:rag & D3\_text\_only & 125 & 110 (88.0\%) & 15 (12.0\%) & 0 (0.0\%) & 0 (0.0\%) & 0 (0.0\%) & 0 (0.0\%) \\
claude-haiku-4-5:rag & D4\_abbreviated & 125 & 95 (76.0\%) & 30 (24.0\%) & 0 (0.0\%) & 0 (0.0\%) & 0 (0.0\%) & 0 (0.0\%) \\
claude-haiku-4-5:zeroshot & D1\_full & 125 & 125 (100.0\%) & 0 (0.0\%) & 0 (0.0\%) & 0 (0.0\%) & 0 (0.0\%) & 0 (0.0\%) \\
claude-haiku-4-5:zeroshot & D2\_no\_code & 125 & 113 (90.4\%) & 8 (6.4\%) & 0 (0.0\%) & 0 (0.0\%) & 0 (0.0\%) & 4 (3.2\%) \\
claude-haiku-4-5:zeroshot & D3\_text\_only & 125 & 105 (84.0\%) & 14 (11.2\%) & 0 (0.0\%) & 0 (0.0\%) & 0 (0.0\%) & 6 (4.8\%) \\
claude-haiku-4-5:zeroshot & D4\_abbreviated & 125 & 102 (81.6\%) & 18 (14.4\%) & 0 (0.0\%) & 1 (0.8\%) & 0 (0.0\%) & 4 (3.2\%) \\
claude-opus-4-7:constrained & D1\_full & 125 & 125 (100.0\%) & 0 (0.0\%) & 0 (0.0\%) & 0 (0.0\%) & 0 (0.0\%) & 0 (0.0\%) \\
claude-opus-4-7:constrained & D2\_no\_code & 125 & 125 (100.0\%) & 0 (0.0\%) & 0 (0.0\%) & 0 (0.0\%) & 0 (0.0\%) & 0 (0.0\%) \\
claude-opus-4-7:constrained & D3\_text\_only & 125 & 116 (92.8\%) & 8 (6.4\%) & 0 (0.0\%) & 0 (0.0\%) & 1 (0.8\%) & 0 (0.0\%) \\
claude-opus-4-7:constrained & D4\_abbreviated & 125 & 108 (86.4\%) & 17 (13.6\%) & 0 (0.0\%) & 0 (0.0\%) & 0 (0.0\%) & 0 (0.0\%) \\
claude-opus-4-7:rag & D1\_full & 125 & 125 (100.0\%) & 0 (0.0\%) & 0 (0.0\%) & 0 (0.0\%) & 0 (0.0\%) & 0 (0.0\%) \\
claude-opus-4-7:rag & D2\_no\_code & 125 & 125 (100.0\%) & 0 (0.0\%) & 0 (0.0\%) & 0 (0.0\%) & 0 (0.0\%) & 0 (0.0\%) \\
claude-opus-4-7:rag & D3\_text\_only & 125 & 109 (87.2\%) & 16 (12.8\%) & 0 (0.0\%) & 0 (0.0\%) & 0 (0.0\%) & 0 (0.0\%) \\
claude-opus-4-7:rag & D4\_abbreviated & 125 & 102 (81.6\%) & 23 (18.4\%) & 0 (0.0\%) & 0 (0.0\%) & 0 (0.0\%) & 0 (0.0\%) \\
claude-opus-4-7:zeroshot & D1\_full & 125 & 124 (99.2\%) & 0 (0.0\%) & 0 (0.0\%) & 0 (0.0\%) & 0 (0.0\%) & 1 (0.8\%) \\
claude-opus-4-7:zeroshot & D2\_no\_code & 125 & 124 (99.2\%) & 1 (0.8\%) & 0 (0.0\%) & 0 (0.0\%) & 0 (0.0\%) & 0 (0.0\%) \\
claude-opus-4-7:zeroshot & D3\_text\_only & 125 & 96 (76.8\%) & 22 (17.6\%) & 0 (0.0\%) & 0 (0.0\%) & 0 (0.0\%) & 7 (5.6\%) \\
claude-opus-4-7:zeroshot & D4\_abbreviated & 125 & 101 (80.8\%) & 17 (13.6\%) & 0 (0.0\%) & 5 (4.0\%) & 2 (1.6\%) & 0 (0.0\%) \\
claude-sonnet-4-6:constrained & D1\_full & 125 & 125 (100.0\%) & 0 (0.0\%) & 0 (0.0\%) & 0 (0.0\%) & 0 (0.0\%) & 0 (0.0\%) \\
claude-sonnet-4-6:constrained & D2\_no\_code & 125 & 125 (100.0\%) & 0 (0.0\%) & 0 (0.0\%) & 0 (0.0\%) & 0 (0.0\%) & 0 (0.0\%) \\
claude-sonnet-4-6:constrained & D3\_text\_only & 125 & 110 (88.0\%) & 15 (12.0\%) & 0 (0.0\%) & 0 (0.0\%) & 0 (0.0\%) & 0 (0.0\%) \\
claude-sonnet-4-6:constrained & D4\_abbreviated & 125 & 113 (90.4\%) & 12 (9.6\%) & 0 (0.0\%) & 0 (0.0\%) & 0 (0.0\%) & 0 (0.0\%) \\
claude-sonnet-4-6:rag & D1\_full & 125 & 125 (100.0\%) & 0 (0.0\%) & 0 (0.0\%) & 0 (0.0\%) & 0 (0.0\%) & 0 (0.0\%) \\
claude-sonnet-4-6:rag & D2\_no\_code & 125 & 125 (100.0\%) & 0 (0.0\%) & 0 (0.0\%) & 0 (0.0\%) & 0 (0.0\%) & 0 (0.0\%) \\
claude-sonnet-4-6:rag & D3\_text\_only & 125 & 108 (86.4\%) & 17 (13.6\%) & 0 (0.0\%) & 0 (0.0\%) & 0 (0.0\%) & 0 (0.0\%) \\
claude-sonnet-4-6:rag & D4\_abbreviated & 125 & 96 (76.8\%) & 29 (23.2\%) & 0 (0.0\%) & 0 (0.0\%) & 0 (0.0\%) & 0 (0.0\%) \\
claude-sonnet-4-6:zeroshot & D1\_full & 125 & 125 (100.0\%) & 0 (0.0\%) & 0 (0.0\%) & 0 (0.0\%) & 0 (0.0\%) & 0 (0.0\%) \\
claude-sonnet-4-6:zeroshot & D2\_no\_code & 125 & 94 (75.2\%) & 30 (24.0\%) & 0 (0.0\%) & 0 (0.0\%) & 0 (0.0\%) & 1 (0.8\%) \\
claude-sonnet-4-6:zeroshot & D3\_text\_only & 125 & 81 (64.8\%) & 36 (28.8\%) & 0 (0.0\%) & 0 (0.0\%) & 0 (0.0\%) & 8 (6.4\%) \\
claude-sonnet-4-6:zeroshot & D4\_abbreviated & 125 & 77 (61.6\%) & 40 (32.0\%) & 0 (0.0\%) & 4 (3.2\%) & 0 (0.0\%) & 4 (3.2\%) \\
gemini-2.5-flash:constrained & D1\_full & 125 & 124 (99.2\%) & 0 (0.0\%) & 0 (0.0\%) & 0 (0.0\%) & 1 (0.8\%) & 0 (0.0\%) \\
gemini-2.5-flash:constrained & D2\_no\_code & 125 & 125 (100.0\%) & 0 (0.0\%) & 0 (0.0\%) & 0 (0.0\%) & 0 (0.0\%) & 0 (0.0\%) \\
gemini-2.5-flash:constrained & D3\_text\_only & 125 & 105 (84.0\%) & 5 (4.0\%) & 0 (0.0\%) & 0 (0.0\%) & 15 (12.0\%) & 0 (0.0\%) \\
gemini-2.5-flash:constrained & D4\_abbreviated & 125 & 88 (70.4\%) & 14 (11.2\%) & 0 (0.0\%) & 0 (0.0\%) & 23 (18.4\%) & 0 (0.0\%) \\
gemini-2.5-flash:rag & D1\_full & 125 & 123 (98.4\%) & 0 (0.0\%) & 0 (0.0\%) & 0 (0.0\%) & 2 (1.6\%) & 0 (0.0\%) \\
gemini-2.5-flash:rag & D2\_no\_code & 125 & 120 (96.0\%) & 0 (0.0\%) & 0 (0.0\%) & 0 (0.0\%) & 5 (4.0\%) & 0 (0.0\%) \\
gemini-2.5-flash:rag & D3\_text\_only & 125 & 96 (76.8\%) & 2 (1.6\%) & 0 (0.0\%) & 0 (0.0\%) & 27 (21.6\%) & 0 (0.0\%) \\
gemini-2.5-flash:rag & D4\_abbreviated & 125 & 84 (67.2\%) & 15 (12.0\%) & 0 (0.0\%) & 0 (0.0\%) & 26 (20.8\%) & 0 (0.0\%) \\
gemini-2.5-flash:zeroshot & D1\_full & 125 & 125 (100.0\%) & 0 (0.0\%) & 0 (0.0\%) & 0 (0.0\%) & 0 (0.0\%) & 0 (0.0\%) \\
gemini-2.5-flash:zeroshot & D2\_no\_code & 125 & 107 (85.6\%) & 4 (3.2\%) & 0 (0.0\%) & 0 (0.0\%) & 14 (11.2\%) & 0 (0.0\%) \\
gemini-2.5-flash:zeroshot & D3\_text\_only & 125 & 76 (60.8\%) & 1 (0.8\%) & 0 (0.0\%) & 0 (0.0\%) & 47 (37.6\%) & 1 (0.8\%) \\
gemini-2.5-flash:zeroshot & D4\_abbreviated & 125 & 69 (55.2\%) & 2 (1.6\%) & 0 (0.0\%) & 0 (0.0\%) & 53 (42.4\%) & 1 (0.8\%) \\
gemini-2.5-pro:constrained & D1\_full & 125 & 57 (45.6\%) & 0 (0.0\%) & 0 (0.0\%) & 0 (0.0\%) & 68 (54.4\%) & 0 (0.0\%) \\
gemini-2.5-pro:constrained & D2\_no\_code & 125 & 85 (68.0\%) & 0 (0.0\%) & 0 (0.0\%) & 0 (0.0\%) & 40 (32.0\%) & 0 (0.0\%) \\
gemini-2.5-pro:constrained & D3\_text\_only & 125 & 40 (32.0\%) & 1 (0.8\%) & 0 (0.0\%) & 0 (0.0\%) & 84 (67.2\%) & 0 (0.0\%) \\
gemini-2.5-pro:constrained & D4\_abbreviated & 125 & 49 (39.2\%) & 0 (0.0\%) & 0 (0.0\%) & 0 (0.0\%) & 76 (60.8\%) & 0 (0.0\%) \\
gemini-2.5-pro:rag & D1\_full & 125 & 49 (39.2\%) & 0 (0.0\%) & 0 (0.0\%) & 0 (0.0\%) & 76 (60.8\%) & 0 (0.0\%) \\
gemini-2.5-pro:rag & D2\_no\_code & 125 & 44 (35.2\%) & 0 (0.0\%) & 0 (0.0\%) & 0 (0.0\%) & 81 (64.8\%) & 0 (0.0\%) \\
gemini-2.5-pro:rag & D3\_text\_only & 125 & 31 (24.8\%) & 0 (0.0\%) & 0 (0.0\%) & 0 (0.0\%) & 94 (75.2\%) & 0 (0.0\%) \\
gemini-2.5-pro:rag & D4\_abbreviated & 125 & 35 (28.0\%) & 7 (5.6\%) & 0 (0.0\%) & 0 (0.0\%) & 83 (66.4\%) & 0 (0.0\%) \\
gemini-2.5-pro:zeroshot & D1\_full & 125 & 28 (22.4\%) & 0 (0.0\%) & 0 (0.0\%) & 0 (0.0\%) & 97 (77.6\%) & 0 (0.0\%) \\
gemini-2.5-pro:zeroshot & D2\_no\_code & 125 & 18 (14.4\%) & 0 (0.0\%) & 0 (0.0\%) & 0 (0.0\%) & 107 (85.6\%) & 0 (0.0\%) \\
gemini-2.5-pro:zeroshot & D3\_text\_only & 125 & 7 (5.6\%) & 0 (0.0\%) & 0 (0.0\%) & 0 (0.0\%) & 118 (94.4\%) & 0 (0.0\%) \\
gemini-2.5-pro:zeroshot & D4\_abbreviated & 125 & 6 (4.8\%) & 0 (0.0\%) & 0 (0.0\%) & 0 (0.0\%) & 119 (95.2\%) & 0 (0.0\%) \\
gemini-3-flash-preview:constrained & D1\_full & 125 & 122 (97.6\%) & 0 (0.0\%) & 0 (0.0\%) & 0 (0.0\%) & 3 (2.4\%) & 0 (0.0\%) \\
gemini-3-flash-preview:constrained & D2\_no\_code & 125 & 121 (96.8\%) & 0 (0.0\%) & 0 (0.0\%) & 0 (0.0\%) & 4 (3.2\%) & 0 (0.0\%) \\
gemini-3-flash-preview:constrained & D3\_text\_only & 125 & 99 (79.2\%) & 4 (3.2\%) & 0 (0.0\%) & 0 (0.0\%) & 22 (17.6\%) & 0 (0.0\%) \\
gemini-3-flash-preview:constrained & D4\_abbreviated & 125 & 106 (84.8\%) & 3 (2.4\%) & 0 (0.0\%) & 0 (0.0\%) & 16 (12.8\%) & 0 (0.0\%) \\
gemini-3-flash-preview:rag & D1\_full & 125 & 119 (95.2\%) & 0 (0.0\%) & 0 (0.0\%) & 0 (0.0\%) & 6 (4.8\%) & 0 (0.0\%) \\
gemini-3-flash-preview:rag & D2\_no\_code & 125 & 107 (85.6\%) & 0 (0.0\%) & 0 (0.0\%) & 0 (0.0\%) & 18 (14.4\%) & 0 (0.0\%) \\
gemini-3-flash-preview:rag & D3\_text\_only & 125 & 102 (81.6\%) & 3 (2.4\%) & 0 (0.0\%) & 0 (0.0\%) & 20 (16.0\%) & 0 (0.0\%) \\
gemini-3-flash-preview:rag & D4\_abbreviated & 125 & 84 (67.2\%) & 17 (13.6\%) & 0 (0.0\%) & 0 (0.0\%) & 24 (19.2\%) & 0 (0.0\%) \\
gemini-3-flash-preview:zeroshot & D1\_full & 125 & 97 (77.6\%) & 0 (0.0\%) & 0 (0.0\%) & 0 (0.0\%) & 28 (22.4\%) & 0 (0.0\%) \\
gemini-3-flash-preview:zeroshot & D2\_no\_code & 125 & 90 (72.0\%) & 0 (0.0\%) & 0 (0.0\%) & 0 (0.0\%) & 35 (28.0\%) & 0 (0.0\%) \\
gemini-3-flash-preview:zeroshot & D3\_text\_only & 125 & 69 (55.2\%) & 4 (3.2\%) & 0 (0.0\%) & 0 (0.0\%) & 52 (41.6\%) & 0 (0.0\%) \\
gemini-3-flash-preview:zeroshot & D4\_abbreviated & 125 & 75 (60.0\%) & 1 (0.8\%) & 0 (0.0\%) & 0 (0.0\%) & 49 (39.2\%) & 0 (0.0\%) \\
gpt-4o-mini:constrained & D1\_full & 125 & 125 (100.0\%) & 0 (0.0\%) & 0 (0.0\%) & 0 (0.0\%) & 0 (0.0\%) & 0 (0.0\%) \\
gpt-4o-mini:constrained & D2\_no\_code & 125 & 125 (100.0\%) & 0 (0.0\%) & 0 (0.0\%) & 0 (0.0\%) & 0 (0.0\%) & 0 (0.0\%) \\
gpt-4o-mini:constrained & D3\_text\_only & 125 & 113 (90.4\%) & 12 (9.6\%) & 0 (0.0\%) & 0 (0.0\%) & 0 (0.0\%) & 0 (0.0\%) \\
gpt-4o-mini:constrained & D4\_abbreviated & 125 & 108 (86.4\%) & 17 (13.6\%) & 0 (0.0\%) & 0 (0.0\%) & 0 (0.0\%) & 0 (0.0\%) \\
gpt-4o-mini:rag & D1\_full & 125 & 125 (100.0\%) & 0 (0.0\%) & 0 (0.0\%) & 0 (0.0\%) & 0 (0.0\%) & 0 (0.0\%) \\
gpt-4o-mini:rag & D2\_no\_code & 125 & 125 (100.0\%) & 0 (0.0\%) & 0 (0.0\%) & 0 (0.0\%) & 0 (0.0\%) & 0 (0.0\%) \\
gpt-4o-mini:rag & D3\_text\_only & 125 & 101 (80.8\%) & 24 (19.2\%) & 0 (0.0\%) & 0 (0.0\%) & 0 (0.0\%) & 0 (0.0\%) \\
gpt-4o-mini:rag & D4\_abbreviated & 125 & 101 (80.8\%) & 24 (19.2\%) & 0 (0.0\%) & 0 (0.0\%) & 0 (0.0\%) & 0 (0.0\%) \\
gpt-4o-mini:zeroshot & D1\_full & 125 & 124 (99.2\%) & 0 (0.0\%) & 0 (0.0\%) & 0 (0.0\%) & 0 (0.0\%) & 1 (0.8\%) \\
gpt-4o-mini:zeroshot & D2\_no\_code & 125 & 106 (84.8\%) & 2 (1.6\%) & 1 (0.8\%) & 0 (0.0\%) & 0 (0.0\%) & 16 (12.8\%) \\
gpt-4o-mini:zeroshot & D3\_text\_only & 125 & 98 (78.4\%) & 13 (10.4\%) & 0 (0.0\%) & 2 (1.6\%) & 0 (0.0\%) & 12 (9.6\%) \\
gpt-4o-mini:zeroshot & D4\_abbreviated & 125 & 98 (78.4\%) & 16 (12.8\%) & 0 (0.0\%) & 2 (1.6\%) & 0 (0.0\%) & 9 (7.2\%) \\
gpt-5.5:constrained & D1\_full & 125 & 125 (100.0\%) & 0 (0.0\%) & 0 (0.0\%) & 0 (0.0\%) & 0 (0.0\%) & 0 (0.0\%) \\
gpt-5.5:constrained & D2\_no\_code & 125 & 125 (100.0\%) & 0 (0.0\%) & 0 (0.0\%) & 0 (0.0\%) & 0 (0.0\%) & 0 (0.0\%) \\
gpt-5.5:constrained & D3\_text\_only & 125 & 109 (87.2\%) & 16 (12.8\%) & 0 (0.0\%) & 0 (0.0\%) & 0 (0.0\%) & 0 (0.0\%) \\
gpt-5.5:constrained & D4\_abbreviated & 125 & 109 (87.2\%) & 16 (12.8\%) & 0 (0.0\%) & 0 (0.0\%) & 0 (0.0\%) & 0 (0.0\%) \\
gpt-5.5:rag & D1\_full & 125 & 124 (99.2\%) & 0 (0.0\%) & 0 (0.0\%) & 0 (0.0\%) & 1 (0.8\%) & 0 (0.0\%) \\
gpt-5.5:rag & D2\_no\_code & 125 & 123 (98.4\%) & 0 (0.0\%) & 0 (0.0\%) & 0 (0.0\%) & 2 (1.6\%) & 0 (0.0\%) \\
gpt-5.5:rag & D3\_text\_only & 125 & 109 (87.2\%) & 16 (12.8\%) & 0 (0.0\%) & 0 (0.0\%) & 0 (0.0\%) & 0 (0.0\%) \\
gpt-5.5:rag & D4\_abbreviated & 125 & 90 (72.0\%) & 29 (23.2\%) & 0 (0.0\%) & 0 (0.0\%) & 6 (4.8\%) & 0 (0.0\%) \\
gpt-5.5:zeroshot & D1\_full & 125 & 124 (99.2\%) & 0 (0.0\%) & 0 (0.0\%) & 0 (0.0\%) & 1 (0.8\%) & 0 (0.0\%) \\
gpt-5.5:zeroshot & D2\_no\_code & 125 & 125 (100.0\%) & 0 (0.0\%) & 0 (0.0\%) & 0 (0.0\%) & 0 (0.0\%) & 0 (0.0\%) \\
gpt-5.5:zeroshot & D3\_text\_only & 125 & 89 (71.2\%) & 16 (12.8\%) & 0 (0.0\%) & 19 (15.2\%) & 1 (0.8\%) & 0 (0.0\%) \\
gpt-5.5:zeroshot & D4\_abbreviated & 125 & 106 (84.8\%) & 16 (12.8\%) & 0 (0.0\%) & 3 (2.4\%) & 0 (0.0\%) & 0 (0.0\%) \\
sentence-transformer:retrieval & D1\_full & 125 & 101 (80.8\%) & 24 (19.2\%) & 0 (0.0\%) & 0 (0.0\%) & 0 (0.0\%) & 0 (0.0\%) \\
sentence-transformer:retrieval & D2\_no\_code & 125 & 101 (80.8\%) & 24 (19.2\%) & 0 (0.0\%) & 0 (0.0\%) & 0 (0.0\%) & 0 (0.0\%) \\
sentence-transformer:retrieval & D3\_text\_only & 125 & 101 (80.8\%) & 24 (19.2\%) & 0 (0.0\%) & 0 (0.0\%) & 0 (0.0\%) & 0 (0.0\%) \\
sentence-transformer:retrieval & D4\_abbreviated & 125 & 42 (33.6\%) & 4 (3.2\%) & 8 (6.4\%) & 65 (52.0\%) & 0 (0.0\%) & 6 (4.8\%) \\
\bottomrule()
\end{longtable}


\subsection{S2.2--S2.7: Pointer to repository}\label{s2.2-s2.7-pointer-to-repository}

Extended per-cell results --- the full hallucination-rate table
with Wilson 95\% CIs (S2.2), the full no\_prediction-rate table
with Wilson 95\% CIs (S2.3), the full 28-configuration top-1
vs.\ top-5 lift table (S2.4), the per-bucket cost decomposition
(S2.5), the per-disease-group breakdown (S2.6 --- deferred to v2
at n=125), and the per-vocabulary-position breakdown (S2.7 ---
similarly deferred) --- are released in the project repository at
\url{https://github.com/ufbfung/phantom-codes} under
\texttt{results/summary/n125\_run\_v2/} alongside the
per-prediction CSV.


\subsection{S2.8: Hallucinated-code listing}\label{s2.8-hallucinated-code-listing}

Every fabricated (non-existent) ICD-10-CM code emitted as a top-1
prediction in the headline run, grouped by (model, mode, code) and
sorted by count. A ``fabricated'' code is one that does not appear in
the FY2026 CMS-published ICD-10-CM tabular list, mechanically
checked by the bundled validator at
\href{https://github.com/ufbfung/phantom-codes/tree/main/src/phantom_codes/data/icd10cm/}{\texttt{src/phantom\_codes/data/icd10cm/}}.
Empty-prediction rows (the model returned no code at all) are
classified as \texttt{no\_prediction} in the 6-way taxonomy and are not
included here; see §S2.3 for those rates.

Across the 28-configuration matrix on n=125 × 4 modes = 14,000
total top-1 predictions, \textbf{352 (2.5\%) were fabrications}, drawn
from 16 unique non-existent ICD-10-CM strings.

\begin{longtable}[]{@{}
  >{\raggedright\arraybackslash}p{(\columnwidth - 8\tabcolsep) * \real{0.1875}}
  >{\raggedright\arraybackslash}p{(\columnwidth - 8\tabcolsep) * \real{0.1875}}
  >{\raggedright\arraybackslash}p{(\columnwidth - 8\tabcolsep) * \real{0.1875}}
  >{\raggedright\arraybackslash}p{(\columnwidth - 8\tabcolsep) * \real{0.1875}}
  >{\raggedleft\arraybackslash}p{(\columnwidth - 8\tabcolsep) * \real{0.2500}}@{}}
\toprule()
\begin{minipage}[b]{\linewidth}\raggedright
Model
\end{minipage} & \begin{minipage}[b]{\linewidth}\raggedright
Mode
\end{minipage} & \begin{minipage}[b]{\linewidth}\raggedright
Predicted code
\end{minipage} & \begin{minipage}[b]{\linewidth}\raggedright
Display
\end{minipage} & \begin{minipage}[b]{\linewidth}\raggedleft
n
\end{minipage} \\
\midrule()
\endhead
baseline:fuzzy & D1\_full & E66.8 & Other obesity & 30 \\
baseline:fuzzy & D2\_no\_code & E66.8 & Other obesity & 30 \\
baseline:fuzzy & D3\_text\_only & E66.8 & Other obesity & 30 \\
baseline:fuzzy & D4\_abbreviated & E66.8 & Other obesity & 30 \\
baseline:tfidf & D1\_full & E66.8 & Other obesity & 30 \\
baseline:tfidf & D2\_no\_code & E66.8 & Other obesity & 30 \\
baseline:tfidf & D3\_text\_only & E66.8 & Other obesity & 30 \\
baseline:tfidf & D4\_abbreviated & E66.8 & Other obesity & 30 \\
baseline:tfidf & D4\_abbreviated & I63.0 & Cerebral infarction due to thrombosis of precerebral arteries & 9 \\
gpt-4o-mini:zeroshot & D2\_no\_code & E11.2 & Type 2 diabetes mellitus with kidney complications & 8 \\
claude-opus-4-7:zeroshot & D3\_text\_only & N18.3 & Chronic kidney disease, stage 3 (moderate) & 7 \\
claude-sonnet-4-6:zeroshot & D3\_text\_only & N18.3 & Chronic kidney disease, stage 3 (moderate) & 7 \\
gpt-4o-mini:zeroshot & D2\_no\_code & N18.3 & Chronic kidney disease, stage 3 (moderate) & 7 \\
baseline:tfidf & D1\_full & E78.4 & Other hyperlipidemia & 6 \\
baseline:tfidf & D2\_no\_code & E78.4 & Other hyperlipidemia & 6 \\
baseline:tfidf & D3\_text\_only & E78.4 & Other hyperlipidemia & 6 \\
gpt-4o-mini:zeroshot & D3\_text\_only & E11.2 & Type 2 diabetes mellitus with kidney complications & 6 \\
sentence-transformer:retrieval & D4\_abbreviated & I73.8 & Other specified peripheral vascular diseases & 6 \\
claude-haiku-4-5:zeroshot & D3\_text\_only & N18.3 & Chronic kidney disease, stage 3 & 5 \\
gpt-4o-mini:zeroshot & D3\_text\_only & N18.3 & Chronic kidney disease stage 3 (moderate) & 5 \\
baseline:fuzzy & D4\_abbreviated & E78.7 & Disorders of bile acid and cholesterol metabolism & 4 \\
claude-haiku-4-5:zeroshot & D2\_no\_code & N18.3 & Chronic kidney disease, stage 3 (moderate) & 3 \\
claude-sonnet-4-6:zeroshot & D4\_abbreviated & N18.3 & Chronic kidney disease, stage 3 (moderate) & 3 \\
gpt-4o-mini:zeroshot & D4\_abbreviated & E11.2 & Type 2 diabetes mellitus with kidney complications & 3 \\
gpt-4o-mini:zeroshot & D4\_abbreviated & N18.3 & Chronic kidney disease stage 3 (moderate) & 3 \\
claude-haiku-4-5:zeroshot & D2\_no\_code & E11.339 & Type 2 diabetes mellitus with mild nonproliferative diabetic retinopathy without macular edema, unspecified eye & 1 \\
claude-haiku-4-5:zeroshot & D3\_text\_only & E11.329 & Type 2 diabetes mellitus with nonproliferative diabetic retinopathy without macular edema & 1 \\
claude-haiku-4-5:zeroshot & D4\_abbreviated & E11.329 & Type 2 diabetes mellitus with mild nonproliferative diabetic retinopathy without macular edema & 1 \\
claude-haiku-4-5:zeroshot & D4\_abbreviated & N18.3 & Chronic kidney disease, stage 3 & 1 \\
claude-haiku-4-5:zeroshot & D4\_abbreviated & Z87.72 & Personal history of diseases of the kidney and ureter & 1 \\
claude-haiku-4-5:zeroshot & D4\_abbreviated & Z87.735 & Personal history of chronic kidney disease & 1 \\
claude-opus-4-7:zeroshot & D1\_full & E11.329 & Type 2 diabetes mellitus with mild nonproliferative diabetic retinopathy without macular edema, unspecified eye & 1 \\
claude-sonnet-4-6:zeroshot & D2\_no\_code & E11.329 & Type 2 diabetes mellitus with mild nonproliferative diabetic retinopathy without macular edema, unspecified eye & 1 \\
claude-sonnet-4-6:zeroshot & D3\_text\_only & E11.32 & Type 2 diabetes mellitus with mild nonproliferative diabetic retinopathy & 1 \\
claude-sonnet-4-6:zeroshot & D4\_abbreviated & E11.32 & Type 2 diabetes mellitus with mild nonproliferative diabetic retinopathy & 1 \\
gemini-2.5-flash:zeroshot & D3\_text\_only & E78.10 & Pure hypertriglyceridemia, unspecified & 1 \\
gemini-2.5-flash:zeroshot & D4\_abbreviated & N18.3 & Chronic kidney disease, stage 3 (moderate) & 1 \\
gpt-4o-mini:zeroshot & D1\_full & E11.359 & Type 2 diabetes mellitus with mild nonproliferative diabetic retinopathy, unspecified eye & 1 \\
gpt-4o-mini:zeroshot & D2\_no\_code & E11.359 & Type 2 diabetes mellitus with mild nonproliferative diabetic retinopathy, unspecified eye & 1 \\
gpt-4o-mini:zeroshot & D3\_text\_only & E11.359 & Type 2 diabetes mellitus with mild nonproliferative diabetic retinopathy, unspecified eye & 1 \\
gpt-4o-mini:zeroshot & D4\_abbreviated & E11.3 & Type 2 diabetes mellitus with renal complications & 1 \\
gpt-4o-mini:zeroshot & D4\_abbreviated & E11.359 & Type 2 diabetes mellitus with nonproliferative diabetic retinopathy, unspecified eye & 1 \\
gpt-4o-mini:zeroshot & D4\_abbreviated & N08.3 & Disorder of kidney due to type 1 diabetes mellitus & 1 \\
\bottomrule()
\end{longtable}

Two qualitative observations from the listing:

\begin{enumerate}
\def\labelenumi{\arabic{enumi}.}
\tightlist
\item
  \textbf{String-matching baselines fabricate by design.} Both
  \texttt{baseline:fuzzy} and \texttt{baseline:tfidf} repeatedly return \texttt{E66.8}
  (``Other obesity'') and similar near-miss codes that exist in
  ICD-10-CM but were not in the bundled FY2026 CMS validator's
  accepted-code set under the strict comparison logic. These are
  fabrications under the \emph{narrow} mechanical definition; whether
  they should be reclassified to a softer ``out-of-domain''
  category in v2 is a methodology choice.
\item
  \textbf{Frontier-model fabrications cluster on plausible but
  non-existent specificity.} Most frontier-model fabrications fall
  into two patterns: (a) confused chronic-kidney-disease subcoding
  (\texttt{N18.3} instead of the real \texttt{N18.30}/\texttt{N18.31}/\texttt{N18.32}); and

  \begin{enumerate}
  \def\labelenumii{(\alph{enumii})}
  \setcounter{enumii}{1}
  \tightlist
  \item
    over-specified type-2-diabetes-with-retinopathy codes
    (\texttt{E11.329}, \texttt{E11.339}, \texttt{E11.359}) that look like legitimate
    ICD-10-CM strings but do not exist in the FY2026 tabular list.
    The pattern is recognizable: the model is reaching for a more
    specific code than the input warrants and inventing a plausible
    suffix.
  \end{enumerate}
\end{enumerate}

\section{S5: Cost economics: pointer to repository}\label{s5-cost-economics-extended-analysis}

The headline per-call cost, cost-per-correct, and Pareto-frontier
analysis appear in main §3. Extended analyses --- the full
per-(model, prompting-mode) cost table, per-bucket cost
decomposition, production-deployment cost projections at multiple
organization scales, the deployment break-even analysis vs. a human
coder, qualitative deployment considerations (Anthropic prompt-
caching threshold, reliability as a cost factor, pricing-snapshot
reproducibility), and the historical pre-headline-run smoke-test
cost data --- are released in the project repository at
\url{https://github.com/ufbfung/phantom-codes} under
\texttt{paper/supp\_extended/} and the per-prediction CSV at
\texttt{results/summary/n125\_run\_v2/cost\_per\_correct.csv}.
Total run cost across the 28 configurations on the n=125 cohort was
\$49.65, dominated by Claude Opus (\$18.18 across its 3 modes) and
the GPT-5.5 / Sonnet 4.6 groupings.

% S3 (reproducibility appendix, trained-model arm details) lifted
% into the companion arXiv technical report — see references.bib
% entry FungPubMedBERT2026.
% S4 (MI-CLAIM checklist) removed — not required by JAMIA.

\bigskip
\hrule
\bigskip

\printbibliography[heading=bibintoc, title={References}]

\end{document}
